\documentclass[12pt]{article}
\usepackage[canadien]{babel}
\usepackage[utf8]{inputenc}
\usepackage[T1]{fontenc}
\usepackage{fancyhdr}
\usepackage{graphicx}
\usepackage{enumerate}
\usepackage{amsmath}
\usepackage{amssymb}
\usepackage{subfigure}
\usepackage{amsmath}
\usepackage{amssymb}
\usepackage{bm}
\usepackage{url}
\usepackage[dvipsnames]{xcolor}
\usepackage{todonotes}
\usepackage{hyperref}

% If your paper is accepted, change the options for the package
% aistats2e as follows:
%
%\usepackage[accepted]{aistats2e}
%
% This option will print headings for the title of your paper and
% headings for the authors names, plus a copyright note at the end of
% the first column of the first page.
\setlength{\parindent}{0cm}
\addtolength{\oddsidemargin}{-2cm}
\addtolength{\evensidemargin}{-2cm}
\setlength{\textwidth}{17.78cm}
\addtolength{\topmargin}{-2.25cm}
\setlength{\textheight}{23.24cm}
\addtolength{\parskip}{5mm}
\pagestyle{fancy}

%************
%* COMMANDS *
%************

\input{../math_commands.tex}

\newif\ifexercise
\exercisetrue
%\exercisefalse

\newif\ifsolution
\solutiontrue
%\solutionfalse

\usepackage{booktabs}
\usepackage[ruled,vlined]{algorithm2e}

\usepackage{amsthm}
\theoremstyle{definition}
\newtheorem{exercise}{Question}%[chapter]
\newtheorem{answer}{Answer} % asterisk to remove ordering

\newcommand{\Exercise}[1]{
\ifexercise#1\fi
}

\definecolor{answer}{rgb}{0.00, 0.12, 0.60}
\newcommand{\Answer}[1]{
\ifsolution\color{answer}\begin{answer}#1\end{answer}\fi
}

\newif\ifexercise
\exercisetrue
%\exercisefalse

\newif\ifsolution
\solutiontrue
%\solutionfalse

\usepackage{enumitem}
\newcommand{\staritem}{
\addtocounter{enumi}{1}
\item[$\phantom{x}^{*}$\theenumi]}
\setlist[enumerate,1]{leftmargin=*, label=\arabic*.}

\begin{document}


\fancyhead{}
\fancyfoot{}

\fancyhead[L]{
  \begin{tabular}[b]{l}
    IFT6135-A2026  \\
    Prof: Aishwarya Agrawal \\
  \end{tabular}
}
\fancyhead[R]{
  \begin{tabular}[b]{r}
    Assignment 2 - Programming  \\
    Recurrent Neural Networks and Transformers  \\
  \end{tabular}
}
\fancyfoot[C]{- Do not distribute -}

\vspace{1cm}

\shorthandoff{:}
{\textbf{Due Date: November 15th, 2026 at 11:00 pm}}\\


\vspace{-0.5cm}
\underline{Instructions}%
\renewcommand{\labelitemi}{\textbullet}

\renewcommand{\labelitemi}{\textbullet}

\begin{itemize}
\item \emph{This assignment is intensive, start well ahead of time!}
\item \emph{For all questions, show your work.}
% Link to hw2_student.ipynb (both questions). To keep this link, upload hw2_student.ipynb as a new version of this
% Drive file (right-click > Manage versions > Upload new version); otherwise paste the new link here.
\item \emph{The practical part is provided as one Jupyter notebook on Google Colab that contains both questions.} \emph{Save a
copy in your own Drive before editing, and run it on a GPU (Runtime $\rightarrow$ Change runtime type $\rightarrow$ T4
GPU). Submit this notebook (File $\rightarrow$ Download $\rightarrow$ Download .ipynb) to Gradescope:
Questions 1 and 2 are both graded from it.}
\item \emph{You must also submit your report (PDF) on Gradescope. Your report must contain answers to Questions 1 and 2. You also need to provide your Google Colab notebook with outputs already generated and public accessibility.}
\item \emph{TAs for this assignment are \textbf{Qichen Huang} and \textbf{Serena Chan}.}
\end{itemize}
\vspace{0.2cm}


\vspace{0.2cm}

\textbf{Link to Notebook: }\href{https://colab.research.google.com/drive/1U3l83ZcKE-iTDYsCHJfDzjuP3okoZihf?usp=sharing}{Link}

% Q1
\begin{exercise}[50] \textbf{Learning to Update Memory with LSTMs}
\Exercise{\label{ex:lstm}

In this question, you will implement an LSTM and investigate whether it can
retain the latest value of a key after repeated updates and irrelevant inputs.
You will use PyTorch. Automatic differentiation, \texttt{nn.Linear},
\texttt{nn.Embedding}, and primitive tensor operations are allowed.
Do not use \texttt{nn.LSTM}, \texttt{nn.LSTMCell}, other built-in recurrent
modules, or a dictionary-based lookup inside your model.

\textbf{Provided files.} Question 1 is the first part of the Assignment 2 notebook (link in the instructions), which
explains the workflow. Complete only the marked TODO blocks of \texttt{lstm.py}, in the notebook cell that starts with
\texttt{\%\%writefile lstm.py}. Do not edit or delete this first line: the autograder reads your code from this cell.
\texttt{memory\_data.py}, \texttt{train.py}, and \texttt{test\_public.py} are
provided. Preserve function signatures, parameter names, and tensor layouts.
Run the \texttt{lstm.py} cell to create the file, then run the public tests and experiment cells in order.

\subsection*{Task and dataset}
There are three keys $A,B,C$ and four possible values $0,1,2,3$.
An update such as $A=2$ is represented by \emph{one} token; keys and values
are not separate successive tokens. A sequence contains six update tokens,
followed by $D$ distractor tokens and one query token. Each distractor is
sampled from four dedicated symbols that carry no update. The target is the
most recent value assigned to the queried key. For example,
\[
 A=1,\ B=0,\ C=2,\ A=3,\ B=1,\ C=0,\ D_0,\ D_2,\ Q_A
 \quad\longrightarrow\quad 3.
\]
The old value $1$ is incorrect, as is $0$, the value of the last update to
another key. The query is presented only at the end of the sequence.

With zero-indexed key $k$ and value $v$, update tokens have ID $1+4k+v$;
query tokens have ID $13+k$; distractor IDs are $16,17,18,19$.
ID $0$ is reserved for right padding, so the vocabulary size is $20$.
Value labels are $0,1,2,3$. The model receives only token IDs and valid
sequence lengths, never the dataset's analysis metadata.

The supplied generator creates 6,000 training and 1,200 validation examples.
Training and validation use $D\in\{0,2,4,8\}$ and
$R\in\{1,2,3\}$ updates to the queried key. Its last update is at
zero-indexed position 2, 3, or 4. Remaining updates refer to other keys;
in particular, the last update token always concerns a different key.
Consecutive updates to the queried key have different values.
The generator balances query-key/target-value pairs and rejects duplicate
token sequences within and across the supplied splits. In training and validation, final update positions are approximately balanced overall, but not guaranteed to be balanced within each $(D,R)$ condition.

\subsection*{Part 1: Implementing the model (30 points)}
For batch size $B$, embedding dimension $E$, and hidden dimension $H$,
the cell input is $x_t\in\mathbb{R}^{B\times E}$ and both states have
shape $B\times H$. Use the gate order $(i,f,g,o)$:
\[
 [a_i,a_f,a_g,a_o]=[x_t,h_{t-1}]W^\top+b,
 \qquad i_t=\sigma(a_i),\ f_t=\sigma(a_f),\ g_t=\tanh(a_g),\ o_t=\sigma(a_o),
\]
\[
 c_t=f_t\odot c_{t-1}+i_t\odot g_t,
 \qquad h_t=o_t\odot\tanh(c_t).
\]
\begin{enumerate}
\item \textbf{LSTM cell (18 points).} Complete \texttt{LSTMCell.forward}:
compute the affine gate inputs in the specified order (6), apply the
appropriate nonlinearities (6), and return the updated hidden and cell states
(6). The affine layer is already instantiated. Your implementation must
support batched inputs and preserve autograd, input dtype, and device.

\item \textbf{Sequence model (12 points).} Complete
\texttt{MemoryLSTM.recurrent} (8) and \texttt{MemoryLSTM.forward} (4).
Reset the hidden and cell states to zero on every forward call.
Process valid timesteps in order. At right-padded positions, preserve
\emph{both} previous states. Return the hidden state after the final valid
token, including the query. Combine the provided embedding, recurrent cell,
and classifier to return four raw logits per example, without a final
softmax. The model uses one recurrent layer, $E=32$ and $H=64$.
\end{enumerate}

Run \texttt{python -m unittest test\_public -v} before training.
Public tests are illustrative; separate tests will check numerical values,
gradients, batch handling, and padding behavior.

\subsection*{Part 2: Training and memory analysis (20 points)}
Use the provided training script and seed 6135. The configuration is
20 epochs, batch size 128, Adam with learning rate 0.003, and gradient-norm
clipping at 1.0. Do not add layers or perform a hyperparameter search.
Select the checkpoint with the lowest validation cross-entropy, keeping
the earlier epoch in a tie. Do not use a test set to select a checkpoint.
CPU execution is supported. The provided evaluation script uses \texttt{model.eval()} and \texttt{torch.no\_grad()} without updating model parameters.

\begin{enumerate}
\item \textbf{Training and evaluation (6 points).} Report training and
validation loss and accuracy curves (2), the selected epoch and selection
rule (2), and the hardware, software versions, total training/validation
time, and model parameter count (2). Use the supplied notebook outputs directly; no additional experiments are required. Save the supplied run outputs.

\item \textbf{Memory analysis (8 points).} Evaluate the same selected
checkpoint, without retraining, on the 15 supplied conditions
$D\in\{0,4,8,16,32\}$ and $R\in\{1,2,3\}$, with 360 examples per
condition. Here the queried key's final update is balanced across positions
2, 3, and 4 in every condition. Thus its position distribution is the same
when comparing different $R$, without exposing one fixed answer position.
Report an accuracy table and plot accuracy
against $D$, with one curve per $R$ (3). For $R\geq2$, also report the
\emph{stale-value rate}: the fraction of all examples in that condition
whose prediction equals the immediately previous value of the queried key
(2). Report this rate as not applicable for $R=1$.
In one short paragraph (about 3--5 sentences), explain the patterns and distinguish observed performance from a guaranteed
memory capability (3). Identify which distractor lengths were absent from
training. Do not assume that every increase in $D$ or $R$ must reduce accuracy. Differences across $R$ should be interpreted descriptively, since earlier update patterns and spacing also vary. These results concern distractors after all updates, not distractors interleaved with updates.

\item \textbf{Failure analysis (6 points).} Show two actual errors when at least two exist
across the test conditions, using the exported example IDs, decoded
sequences, true labels, predictions, and previous values (3). Prefer different error types when available; both types are not required.
Distinguish stale-value errors from other errors and give a plausible
explanation connected to retaining and updating memory in 2--3 sentences per example (3). You do not need to prove the internal cause.
If fewer than two errors exist, show all available errors and state their
number; if none exist, show two correct examples and discuss the limits of
this evidence. Failure to make errors does not lose marks.
\end{enumerate}

\textbf{Submission.} Submit the notebook (.ipynb) on Gradescope: Part 1 is autograded from the
\texttt{\%\%writefile lstm.py} cell. Answer Part 2 in your PDF report, including the plots, results table and
examples produced by the notebook.
Experimental work is graded on correctness, reproducibility, and
interpretation, not on reaching a prescribed accuracy.
}
\end{exercise}


% Q2
\begin{exercise}[50] \textbf{Word-Level vs.\ Sentence-Level Representation Learning with Transformers}
\Exercise{\label{ex:transformer}}
\sloppy

\textbf{Task:} In this question, you will implement a bidirectional Transformer encoder from its equations using
PyTorch, and train three copies of it on the same sentences with three different objectives: masked language
modeling (MLM), natural language inference (NLI) classification, and NLI contrastive learning. You will then compare
what each objective teaches the model, at the level of single words and at the level of whole sentences.

\textbf{Submission:} Question 2 is the second part of the Assignment 2 notebook (link in the instructions). Complete only the
\texttt{\# TODO} blocks inside the cells whose first line is \texttt{\# GRADED CELL}; everything else is provided. Do
not edit or delete the \texttt{\# GRADED CELL} lines, and do not change any function names, parameter names or tensor
layouts. On Gradescope, submit the notebook (the same file as for Question 1), which is autograded for Parts 1 and 2,
and answer Part 3 in your PDF report.

\textbf{Dataset.}

All three models are trained on the AllNLI corpus (SNLI + MultiNLI): 1{,}147{,}270 unique English sentences of at most
64 tokens (12.6 on average). Each sentence is lowercased, split into words, and every word is replaced by its id in a
vocabulary of the $V = 20{,}000$ most frequent words. Ids 0 to 3 are reserved for \texttt{[PAD]}, \texttt{[MASK]},
\texttt{[CLS]} and \texttt{[UNK]} (words outside the vocabulary), and every sentence starts with \texttt{[CLS]}. The
three objectives use the same sentences in different forms:
\begin{itemize}
  \item \textbf{MLM} uses the sentences on their own, without labels.
  \item \textbf{NLI classification} uses 942{,}069 (premise, hypothesis, label) pairs, where the label says whether the
        hypothesis follows from the premise: \emph{entailment}, \emph{neutral} or \emph{contradiction}.
  \item \textbf{NLI contrastive} uses 557{,}850 (anchor, positive, negative) triplets: a premise, a hypothesis it
        entails, and a hypothesis it contradicts.
\end{itemize}
For evaluation, we use 5{,}000 held-out AllNLI sentences that no model was trained on, and the STS-B benchmark:
1{,}379 pairs of sentences, each rated by humans from 0 (different meaning) to 5 (same meaning).

\textbf{Notation.}

A batch holds $B$ sentences padded to the same length $T$. $x \in \{0, \dots, V-1\}^{B \times T}$ are the token ids
and $a \in \{0, 1\}^{B \times T}$ the attention mask ($a_{b,t} = 1$ at real tokens, $0$ at \texttt{[PAD]}). The
encoder has hidden size $D = 128$, $H = 4$ attention heads of size $d = D/H = 32$, and $L = 4$ blocks. It maps each
sentence to one hidden vector per token, $O = \mathrm{Enc}(x) \in \mathbb{R}^{T \times D}$ with rows
$o_1, \dots, o_T$. The three objectives use this output differently:
\begin{itemize}
  \item \textbf{MLM} hides some words and predicts each one from its context, from the hidden vector at its position
        (one prediction per word).
  \item \textbf{NLI classification} averages the hidden vectors of each sentence into one sentence embedding (mean
        pooling),
        \[
          z = \frac{\sum_{t=1}^{T} a_t\, o_t}{\sum_{t=1}^{T} a_t} \in \mathbb{R}^{D},
        \]
        and predicts the NLI label of a pair from the two embeddings, through a small classifier (one prediction per
        pair of sentences).
  \item \textbf{NLI contrastive} uses the same sentence embeddings, and trains each sentence to be closer, in cosine
        similarity, to the sentence it entails than to all the other sentences of the batch (one prediction per
        sentence).
\end{itemize}
The cosine similarity of two vectors $u, v \in \mathbb{R}^D$ is
$\cos(u, v) = \dfrac{u^\top v}{\|u\|\,\|v\|} \in [-1, 1]$: it is 1 when $u$ and $v$ point the same way, 0 when they
are perpendicular (aka unrelated), and $-1$ when they point in opposite directions.

\textbf{Part 1: Implementing the model (30 pts).}

\begin{enumerate}
  \item \textbf{Batching and padding (2 pts).} Implement \texttt{collate(seqs, device)}. Given $B$ token-id sequences
  $s_1, \dots, s_B$ of lengths $n_1, \dots, n_B$, with $T = \max_b n_b$, return
  \[
    x_{b,t} = \begin{cases} s_{b,t} & t \le n_b \\ 0 \ (\texttt{[PAD]}) & t > n_b \end{cases}, \qquad
    a_{b,t} = \mathbf{1}[t \le n_b].
  \]

  \item \textbf{Embeddings (4 pts).} Implement \texttt{sinusoidal\_positions(T, D)}, which returns the fixed
  positional encodings $PE \in \mathbb{R}^{T \times D}$:
  \[
    PE_{pos,\,2i} = \sin\!\left(\frac{pos}{10000^{2i/D}}\right), \qquad
    PE_{pos,\,2i+1} = \cos\!\left(\frac{pos}{10000^{2i/D}}\right),
  \]
  for positions $pos = 0, \dots, T-1$ and $i = 0, \dots, \tfrac{D}{2}-1$.
  Then implement \texttt{Embeddings.forward(ids)}, the input of the first block. With $E \in \mathbb{R}^{V \times D}$
  the token-embedding matrix and $E_{x_t}$ its row for token $x_t$,
  \[
    X_t = \mathrm{Dropout}\big(\sqrt{D}\, E_{x_t} + PE_t\big), \qquad X \in \mathbb{R}^{T \times D}.
  \]

  \item \textbf{Multi-head attention and the bidirectional mask (7 pts).} Implement
  \texttt{MultiHeadAttention.forward(x, mask)}. Each head
  $h = 1, \dots, H$ has its own query, key and value projections, given to you as lists of \texttt{nn.Linear(D, d)}
  layers \texttt{W\_q[h]}, \texttt{W\_k[h]}, \texttt{W\_v[h]}, and the output projection
  \texttt{W\_o} is an \texttt{nn.Linear(D, D)}:
  \[
    Q_h = W_Q^{(h)}(X), \quad K_h = W_K^{(h)}(X), \quad V_h = W_V^{(h)}(X) \in \mathbb{R}^{T \times d},
  \]
  \[
    A_h = \mathrm{Dropout}\!\left(\mathrm{softmax}\!\left(\frac{Q_h K_h^\top}{\sqrt{d}} + M\right)\right) V_h
    \in \mathbb{R}^{T \times d},
  \]
  \[
    \mathrm{MHA}(X) = W_O\big([A_1, \dots, A_H]\big) \in \mathbb{R}^{T \times D},
  \]
  where the softmax is taken over the keys (each row sums to 1), $[A_1, \dots, A_H]$ concatenates the heads along
  the feature dimension, and $M \in \mathbb{R}^{T \times T}$ is the attention mask: $M_{t,s} = 0$ if query $t$ may
  attend to key $s$, and $M_{t,s} = -10^4$ otherwise (its softmax weight becomes $\approx 0$).


  Our encoder is bidirectional: every
  token may attend to every other token, before or after it. The only restriction is padding, which must never be
  attended to. Implement
  \texttt{bidirectional\_mask(attention\_mask)}, which returns the boolean mask
  $m \in \{0, 1\}^{B \times T \times T}$ with $m_{b,t,s} = a_{b,s}$ ($M_{t,s} = 0$ where $m_{b,t,s} = 1$).

  \item \textbf{Encoder block and encoder (8 pts).} Implement \texttt{EncoderBlock.forward(x, mask)}, a post-norm Transformer
  block:
  \[
    X \leftarrow \mathrm{LN}_1\big(X + \mathrm{Dropout}(\mathrm{MHA}(X))\big), \qquad
    X \leftarrow \mathrm{LN}_2\big(X + \mathrm{Dropout}(\mathrm{FFN}(X))\big),
  \]
  where the feed-forward sublayer $\mathrm{FFN}(X) = W_2\, \mathrm{GELU}(W_1 X)$, with hidden size 512, is provided.


  Then implement \texttt{Encoder.forward(ids, attention\_mask, return\_logits)}: build
  the mask, embed the ids, apply the $L$ blocks in order to obtain $O$, then the MLM head. Thus, the score of word $w$ at position $t$ is
  \[
    \ell_{t,w} = o_t^\top E_w + c_w, \qquad \ell = O E^\top + c \in \mathbb{R}^{T \times V},
  \]
  where $c \in \mathbb{R}^V$ is a learned bias. A word gets a high score when the hidden vector is close to that
  word's embedding.

  \item \textbf{Masked language modeling (9 pts).} Implement \texttt{mask\_tokens(ids, attention\_mask, vocab\_size, p)} as follows:
  \begin{enumerate}
    \item Select each real token independently with probability $p = 0.15$, never \texttt{[PAD]} or a special token
          (ids 0 to 3, e.g.\ \texttt{[CLS]}). Let $\mathcal{M}$ be the set of selected positions $(b, t)$.
    \item For each selected position, independently: with probability $0.8$ replace it with \texttt{[MASK]}; with
          probability $0.1$ replace it with a word id drawn uniformly from $\{4, \dots, V-1\}$; with probability
          $0.1$ leave it unchanged.
    \item Return the corrupted ids $\tilde{x}$ and the labels $y$, where $y_{b,t} = x_{b,t}$ at selected positions
          and $y_{b,t} = -100$ (ignored) elsewhere.
  \end{enumerate}


  Then implement \texttt{cross\_entropy(logits, targets, ignore\_index)} by hand
  (all three objectives use it). For scores $z \in \mathbb{R}^{N \times C}$ and targets $y \in \{0, \dots, C-1\}^N$,
  averaged over the rows $\mathcal{K} = \{n : y_n \ne -100\}$:
  \[
    \mathcal{L}_{\text{CE}}(z, y) = -\frac{1}{|\mathcal{K}|} \sum_{n \in \mathcal{K}}
      \left( z_{n, y_n} - \log \sum_{c=1}^{C} e^{z_{n,c}} \right).
  \]
  Use \texttt{torch.logsumexp} for the log-sum. Then implement
  \texttt{mlm\_loss(logits, labels)}, the cross-entropy over the masked positions only. With logits
  $\ell \in \mathbb{R}^{B \times T \times V}$ computed from the corrupted ids $\tilde{x}$:
  \[
    \mathcal{L}_{\text{MLM}} = -\frac{1}{|\mathcal{M}|} \sum_{(b,t) \in \mathcal{M}}
      \log \mathrm{softmax}(\ell_{b,t})_{\,y_{b,t}}.
  \]

  \item \textbf{Sentence-level objectives (provided).} Read and understand the provided functions, which you will use in
  Part 2:
  \begin{itemize}
    \item \texttt{mean\_pool} and \texttt{sentence\_embedding}: the sentence embedding $z$ defined above.
    \item \texttt{pair\_features(u, v)}: the input of the NLI classifier for a premise embedding $u$ and a
          hypothesis embedding $v$.
    \item \texttt{info\_nce(z\_a, z\_pos, z\_neg, tau)}: the InfoNCE loss with in-batch and hard negatives. For
          anchors $z^a_i$, positives $z^+_i$ (entailments) and hard negatives $z^-_i$ (contradictions),
          $i = 1, \dots, B$, and temperature $\tau = 0.05$:
          \[
            \mathcal{L}_{\text{con}} = -\frac{1}{B} \sum_{i=1}^{B} \log
              \frac{\exp\!\big(\cos(z^a_i, z^+_i)/\tau\big)}
                   {\sum_{j=1}^{B} \exp\!\big(\cos(z^a_i, z^+_j)/\tau\big)
                    + \sum_{j=1}^{B} \exp\!\big(\cos(z^a_i, z^-_j)/\tau\big)}.
          \]
  \end{itemize}
\end{enumerate}

\textbf{Part 2: Losses and training (6 pts).}

The training loop (optimizer, logging, validation) is provided. Implement:

\begin{enumerate}
  \item \textbf{Training step (4 pts).} Implement \texttt{compute\_loss(model, batch, task, vocab\_size, tau)}, which
  returns the loss of one batch:
  \begin{itemize}
    \item \texttt{task = "mlm"}: mask the batch with \texttt{mask\_tokens}, encode the corrupted ids, and
          return $\mathcal{L}_{\text{MLM}}$.
    \item \texttt{task = "nli\_cls"}: compute the sentence embeddings $u$ of the premises and $v$ of the hypotheses,
          and return $\mathcal{L}_{\text{CE}}\big(\texttt{nli\_head}([u;\, v;\, |u - v|]),\ y\big)$, with labels
          $y \in \{0, 1, 2\}$ (entailment, neutral, contradiction).
    \item \texttt{task = "contrastive"}: compute the sentence embeddings of the anchors, positives and negatives,
          and return $\mathcal{L}_{\text{con}}$.
  \end{itemize}

  \item \textbf{Learning-rate schedule (2 pts).} Implement \texttt{lr\_schedule(step, total\_steps, warmup\_steps)},
  a linear warmup followed by a cosine decay. At step $s = 0, 1, \dots, S - 1$, with $S$ total steps and $W$ warmup
  steps, the learning rate is $\eta \cdot m(s)$, where $\eta$ is the peak learning rate and
  \[
    m(s) = \begin{cases}
      \dfrac{s + 1}{W} & \text{if } s < W, \\[2ex]
      \dfrac{1}{2}\left(1 + \cos\!\left(\pi\, \dfrac{s - W}{S - W}\right)\right) & \text{otherwise.}
    \end{cases}
  \]
\end{enumerate}

Then train the three models with the configuration given in the notebook (about 5 minutes each on a T4): each one
starts from a random initialization and uses AdamW (peak learning rate $\eta = 5 \times 10^{-4}$, weight decay $0.01$,
batch size 128, $W = 5\%$ of the steps, gradient clipping at norm 1.0).
\begin{itemize}
  \item \textbf{MLM:} trained with $\mathcal{L}_{\text{MLM}}$ for 6{,}000 steps.
  \item \textbf{NLI classification:} trained with $\mathcal{L}_{\text{CE}}$ on the NLI labels for 5{,}000 steps.
  \item \textbf{NLI contrastive:} trained with $\mathcal{L}_{\text{con}}$ ($\tau = 0.05$) for 4{,}000 steps.
\end{itemize}
An untrained encoder with the same architecture (random initialization) is kept as a baseline.

\textbf{Part 3: Evaluation (14 pts).} Answer in your PDF report. All metrics, tables and plots are computed by
provided functions in the notebook. Include the outputs each question asks for.

The notebook compares the four models with two metrics:
\begin{itemize}
  \item \textbf{Masked-word accuracy}: 15\% of the words of the 5{,}000 held-out sentences are
        masked as in Part 1, and the accuracy is the fraction of masked positions where the MLM head's top guess,
        $\arg\max_w \ell_{t,w}$, is the original word.
  \item \textbf{Cosine of similar vs.\ dissimilar pairs}: let $\mathcal{S}$ be the 338 STS-B
        pairs rated 4 or more (similar) and $\mathcal{D}$ the 308 pairs rated 1 or less (dissimilar). For a pair of
        sentences with embeddings $(z, z')$, the notebook reports the mean cosines and their gap,
        \[
          \overline{\cos}_{\mathcal{S}} = \frac{1}{|\mathcal{S}|} \sum_{(z, z') \in \mathcal{S}} \cos(z, z'), \qquad
          \overline{\cos}_{\mathcal{D}} = \frac{1}{|\mathcal{D}|} \sum_{(z, z') \in \mathcal{D}} \cos(z, z'), \qquad
          \text{gap} = \overline{\cos}_{\mathcal{S}} - \overline{\cos}_{\mathcal{D}}.
        \]
        A good sentence embedding gives similar pairs a much higher cosine than dissimilar ones: a large gap.
\end{itemize}

\begin{enumerate}
  \item \textbf{Training (2 pts).} Include the loss curves (training and validation loss of each model against
  training steps).
  \begin{enumerate}
    \item Did any of the three models overfit? Justify your answer using the curves.
  \end{enumerate}

  \item \textbf{Word level vs.\ sentence level (7 pts).} Include the results table, the masked-word accuracy plot,
  the cosine plot, and the fill-in-the-blank output for the two given sentences and one sentence of your own.
  \begin{enumerate}
    \item Which model has the highest masked-word accuracy? Explain. [2 sentences max.]
    \item Which has the largest cosine gap? Explain. [2 sentences max.]
  \end{enumerate}

  \item \textbf{Classification vs.\ contrastive (3 pts).} Include the search output for the two given queries and
  one query of your own.
  \begin{enumerate}
    \item Both NLI models were trained on the same sentences with the same labels. Explain the differences in cosine gap and search results. [2 sentences max.]
  \end{enumerate}

  \item \textbf{Pretraining and fine-tuning (2 pts).} In practice, sentence encoders are first pretrained with MLM
  and then fine-tuned with a contrastive loss. Based on your results and on the data each objective needs, explain
  what each stage contributes and why they are applied in this order. [2 sentences max.]
\end{enumerate}

\end{exercise}

\end{document}
